\phantomsection
\addcontentsline{toc}{chapter}{\bf ABSTRACT}

\begin{center}
    \large{\bf ABSTRACT}
\end{center}

Optimizing customer conversion paths is a central problem in computational marketing, where user interactions are modeled as a directed graph and the objective is to identify the highest-probability route from entry to conversion. As graph scale increases, exhaustive search becomes intractable. This work proposes a principled pruning framework for log-probability graphs, providing a formal characterization of the trade-off between search efficiency and path reachability.

A \textbf{Probability-Pruned Dijkstra} variant is proposed that introduces a threshold parameter $\tau$ to discard low-probability partial paths during traversal, reducing the effective edge set from $|E|$ to $|E'|\leq|E|$ with guaranteed convergence to the baseline as $\tau\to 0$. Five formal correctness theorems are established, including an all-or-nothing property: the pruned algorithm either returns the globally optimal path or reports no path at all.

Evaluation spans 3,000 configurations---2,160 synthetic (graphs up to 10,000 nodes) and 840 real-data rows from RetailRocket and RecSys~2015---under fixed and adaptive $\tau$ regimes. Across all 706 path-found rows the optimality gap is $0.0000\%$, consistent with the all-or-nothing property (Theorem~4.5). When the pruned algorithm successfully recovers a path, wall-clock speedups range from $1.6\times$ to $7.2\times$ on synthetic data and up to $5.8\times$ on real-world graphs. When the threshold prunes the optimal path entirely, the algorithm detects infeasibility orders of magnitude faster than the baseline (up to $738\times$ on synthetic data and $18{,}882\times$ on real-world item-level graphs). On large sparse real-data graphs, an adaptive-$\tau$ regime raises path-found rates from near~0\% to over~80\%, demonstrating that threshold-based pruning provides both fast infeasibility detection and correct path recovery with guaranteed optimality when a path is returned.
